% LSPV-4 REPORT::FIBERDICH1::2026-10-12T02:36:00Z
% !TEX program = xelatex
\documentclass[11pt]{article}
\usepackage[a4paper,margin=0.88in]{geometry}
\usepackage{amsmath,amssymb,amsthm,booktabs,array,enumitem,xurl,hyperref,microtype}
\hypersetup{colorlinks=true,linkcolor=blue,urlcolor=blue,citecolor=blue}
\newtheorem{definition}{Definition}
\newtheorem{theorem}{Theorem}
\newtheorem{corollary}{Corollary}
\newcommand{\SAT}{\mathsf{SAT}}
\newcommand{\EXACT}{\textbf{PROVED/EXACT}}
\newcommand{\KNOWN}{\textbf{KNOWN THEOREM}}
\newcommand{\COMP}{\textbf{COMPUTATIONAL}}
\newcommand{\COND}{\textbf{CONDITIONAL}}
\newcommand{\CAND}{\textbf{CANDIDATE}}
\newcommand{\OPEN}{\textbf{OPEN}}
\newcommand{\RETIRED}{\textbf{FALSIFIED/RETIRED}}
\title{FIBERDICH1: Post-source challenge obligations either collapse to the decision bit or exceed SAT decision semantics}
\author{LSPV-4 --- Universal Core Red Team}
\date{12 October 2026 (UTC)}

\begin{document}
\maketitle

\begin{center}
\fbox{\parbox{0.94\linewidth}{
\textbf{Executive status.} P versus NP remains open.  This run formalizes the
exact boundary facing the proposed repair in which a source disappears before
a later challenge.  For an arbitrary total response relation, a responder
retaining only the SAT decision bit exists if and only if every challenge has
a common valid response across each same-answer source fiber.  If that
intersection condition holds, one bit defeats every superconstant retention
claim.  If it fails, the protocol asks for source information that the SAT
decision function does not determine; an external compiler may retain such
information, but this is compiler work rather than state forced on every
one-bit SAT decider.  Thus source deletion creates a valid restricted online
problem but does not by itself provide the missing static-SAT transfer.}}
\end{center}

\section{Claim card and exact quantifiers}

\begin{description}[leftmargin=2.4cm,style=nextline]
\item[Claim ID.] \texttt{DECISION-FIBER-ONLINE-DICHOTOMY1}.
\item[Status before run.] \CAND: deleting the source before a later challenge
may force a universal SAT decider to preserve a compatibility/pointer object.
\item[Input/model.] A deterministic one-way cut model.  A finite or arbitrary
source $x\in X$ is available before the cut and unavailable afterward; the
post-cut machine receives challenge $c\in C$ and returns $y\in Y$.
\item[Decision semantics.] A fixed bit $b:X\to\{0,1\}$, instantiated below by
$b(F)=\mathbf 1[F\in\SAT]$.
\item[Response relation.] For every $(x,c)$, a nonempty allowed set
$R(x,c)\subseteq Y$.
\item[Exact prefix.] For all $X,C,Y,b,R$ satisfying totality, characterize
whether there exists $d:\{0,1\}\times C\to Y$ such that
\[
 \forall x\in X\ \forall c\in C:\quad d(b(x),c)\in R(x,c).
 \tag{1}
\]
\item[SAME\_REALIZATION.] Source $x$, decision bit $b(x)$, retained cut state,
challenge $c$, response $y$, and every acquisition/validation/replay cost must
belong to the same declared protocol.  The theorem does not charge an external
compiler's state to the original SAT execution.
\item[Weakest useful result.] An if-and-only-if test separating answer-derived
obligations from source-specific obligations.
\end{description}

\section{Exact decision-fiber theorem}

For $a\in\{0,1\}$ let $X_a=\{x\in X:b(x)=a\}$.  Empty fibers may be omitted.

\begin{theorem}[\texttt{DECISION-FIBER-ONLINE-DICHOTOMY1} --- \EXACT]
There exists a deterministic answer-only responder satisfying (1) if and only
if
\[
 \forall a\in b(X)\ \forall c\in C:\qquad
 \bigcap_{x\in X_a}R(x,c)\ne\varnothing.
 \tag{2}
\]
When (2) holds, the only source-dependent retained state needed by this
responder is $b(x)$, one bit.  When (2) fails, no post-cut responder whose
source-dependent state is only $b(x)$ can be correct for every source and
challenge.
\end{theorem}

\begin{proof}
If (2) holds, choose for each $(a,c)$ one element
$d(a,c)\in\bigcap_{x\in X_a}R(x,c)$.  Then for any $x$ with $b(x)=a$, the
chosen response lies in $R(x,c)$, proving (1).  Conversely, if $d$ satisfies
(1), then for fixed $(a,c)$ the single value $d(a,c)$ lies in every
$R(x,c)$ for $x\in X_a$, so their intersection is nonempty.  If (2) fails,
this converse gives an immediate contradiction to the existence of such a
$d$.  The implementation stores $b(x)$ and later evaluates the fixed table or
algorithm $d$ on $(b(x),c)$; hence it retains one source-dependent bit.
\end{proof}

The theorem is set-theoretic and does not assert that an arbitrary infinite
choice table is efficiently computable.  In every complexity application the
representation and evaluation cost of $d$ must be declared.  This caveat does
not weaken the negative transfer: if a polynomial protocol claims that its
response is forced by the SAT answer alone, it must exhibit an efficient
answer-fiber choice.  If no common response exists, the relation is not a
function of SAT decision semantics.

\begin{corollary}[Static-SAT transfer boundary]
Let $X$ be a family of encoded CNFs and $b(F)=\SAT(F)$.  A post-source
challenge relation falls into exactly one of the following classes.
\begin{enumerate}[label=(\roman*),leftmargin=2em]
\item If (2) holds with an efficiently evaluable choice, a one-bit
predecision summary answers every challenge; any claimed
$\omega(1)$ source-state obligation is false for that relation.
\item If (2) fails, the relation distinguishes at least two same-answer
formulas.  Correct response therefore needs source-specific information
beyond the decision bit.  The need may yield a valid communication or
streaming lower bound in the declared cut model, but it is not implied by the
external behavior of an arbitrary exact SAT decider.
\end{enumerate}
\end{corollary}

In case (ii), an external source compiler can of course inspect $F$ and retain
a message.  That construction, message size, preprocessing, challenge
interface, response computation, and endpoint must all be charged to the new
protocol.  Calling this message ``solver state'' would violate
SAME\_REALIZATION.  A one-bit implementation, recognizable-island shortcut,
or extensionally equivalent wrapper remains free to discard it.

\section{Application to the current online repair}

TESTS-0019 proves that a unique bit-local pointer can be validated and erased;
MODELBRIDGE-0020 proves that auxiliary pointer metadata can be projected away.
The proposed \texttt{SOURCE-SEMANTIC-ONLINE-OBLIGATION1} tries to repair those
failures by making the source unavailable before a causally later challenge.
The theorem above gives the required first audit:
\begin{itemize}[leftmargin=1.6em]
\item If the response sets have common choices on both SAT-answer fibers, the
repair collapses to one bit and \texttt{ANSWER-DERIVED-ONLINE-STATE1} is
\RETIRED as a superconstant universal obligation.
\item If a same-answer fiber has empty intersection for some challenge, the
obligation is genuinely source-specific.  This avoids the one-bit collapse,
but the protocol has added a function not determined by SAT membership.  It
still needs a model-preserving compiler forcing every static SAT computation
to realize that function.
\end{itemize}

Under the temporary premise $P=NP$, if the source-specific response relation
is polynomially balanced, polynomial-time verifiable, and total, STATICNP1
applies: NP prefix search produces a response in polynomial time.  The
polynomial exponent may depend on the assumed SAT decider, so this is no
FLMV contradiction.  If response evaluation is itself placed outside
polynomial time by premise, the desired conclusion has been moved into the
protocol unless an independent, same-model lower bound and static compiler
are proved.

The smallest honest repair is therefore
\texttt{SOURCE-FIBER-SEPARATING-COMPILER1}: give an explicit full-language
relation failing (2), prove a bit-local polynomial compiler from ordinary
static 3SAT whose same realization must answer it after source removal, charge
all state/reacquisition/recomputation, and land that exact machine in an
unconditional forbidden region.  This target is \OPEN and
separation-strength.

\section{Reproducible hostile audit}

The dependency-free program
\begin{verbatim}
timeout 30s python3 experiments/review/decision_fiber_online_dichotomy_audit.py
\end{verbatim}
exhausted every total nonempty binary response relation on four sources with
answer fibers $(2,2)$, two challenges, and two replies.  There are
$3^8=6{,}561$ such relations.  For each relation it compared (2) against all
16 answer-only responders, checked the answer partition directly, and tested
all 15 source partitions.  It made $111{,}537$ assertions with zero failures:
$2{,}401$ relations factored through the answer fibers and $4{,}160$ were
source-specific.  Minimum required message-class counts were
\[
 1:961,qquad 2:4{,}424,qquad 3:1{,}152,qquad 4:24.
\]

The second battery used satisfiable and unsatisfiable Horn, 2SAT,
affine/Tseitin, and bounded-treewidth sources; 48 DIRECT/COPY/DEAD/DELAY/
INTERP/RECOMP encoding variants; 64 independent products; redundant-padding
flags; and 144 one-bit outputs.  Answer-only and universal response relations
passed (2).  Source-bit, wrapper-trace, and encoding-tag relations failed (2)
while remaining trivial to compute from their added metadata.  This is the
hostile point: requiring information beyond the SAT bit does not make that
information computationally hard.

The audit is \COMP and finite.  It proves no asymptotic runtime lower bound.
Script SHA-256:
\begin{verbatim}
7a701d9c1424d903fe516839d159d754598e23673dcb97d21476836e74f1af81
\end{verbatim}
Result SHA-256:
\begin{verbatim}
f437b0bc4a0ff291ba3d42959be4e51aea4a2c8b12176261c80a7bc0ad69011e
\end{verbatim}

\section{Literature and transfer classification}

Yao's communication model explicitly fixes which party sees each input and
what message crosses the cut \cite{Yao79}.  Streaming lower bounds similarly
depend on a declared arrival/access and memory model \cite{BKM18}.  These
theories transfer exactly to a source-removal protocol only after the source,
challenge, message, error criterion, and access rules are fixed.  They do not
themselves prove that an unrestricted random-access SAT decider implements
that protocol.

\begin{itemize}[leftmargin=1.6em]
\item \textbf{Directly transferable theorem.}  The response-fiber
intersection criterion is exactly the deterministic one-way feasibility
condition for an answer-only message.  Yao-style message lower bounds apply
to an explicitly compiled source/challenge function in their declared model.
\item \textbf{Useful analogy requiring a new transfer theorem.}  Control
observability, streaming memory, distributed cuts, VLSI routing, and
cell-probe communication all ask which pre-cut states remain distinguishable.
The missing correspondence is a compiler proving that every ordinary static
SAT computation must preserve the same source partition after the same cut.
\item \textbf{Misleading analogy / failed transfer.}  Physical locality,
network routing, no-reread memory, cryptographic binding, random oracles,
proof-system size, and branching-program lower bounds cannot be imported into
unrestricted deterministic SAT without a model-preserving simulation.
\end{itemize}

The PQ lossiness ledger applies exactly to a declared message map
$L:X\to M$: collision fibers measure which sources share a state.  But when
(2) holds, the SAT bit is already a sufficient two-state quotient; when it
fails, collision entropy describes the new source-response problem, not SAT
decision time.  Special-soundness collision extraction remains a failed
transfer from one deterministic decision transcript.  Schur squares,
quadratic relations, syzygies, and Hilbert functions are unrelated unless a
forced algebraic representation is first proved; hidden structure may instead
give a shortcut.

The theorem and its counterexamples relativize and algebraize.  They define
no large useful constructive circuit property, so Natural Proofs is not the
hinge \cite{RR97}.  Baker--Gill--Solovay and Aaronson--Wigderson therefore
remain active barriers to promoting this semantics-only dichotomy into a
separation \cite{BGS75,AW09}.  FLMV can be invoked only after one ordinary
machine, with the source compiler and all online costs included, satisfies its
simultaneous time--space hypotheses \cite{FLMV05}.

\section{CORPUS REVIEWED}

The run surveyed the complete GitHub tree at pre-run head
\texttt{da8578ce458508ac5d00d2535d783b6beea44929}; read root
\texttt{main.tex}, actual \texttt{main4.tex} and \texttt{main5.tex},
\texttt{AGENTS.md}, \texttt{CODEX.md}, the PNP research skill and protocols,
and all current task specifications.  It read the newest verified MAIN and
SYNTHESIS-0019, newest reports LB-0022, TESTS-0019, PADQUOT1, FORKFACTOR1,
and MODELBRIDGE-0020, and verified that predecessor PADQUOT1 is substantive
with its expected header and final sentinel.  Material historical review
included REREAD1, REACQ1, STATICNP1, ORDERBIND1, TAGGLUE1, the HEC/VSC-HEC/
G-PPC-HEC/CRS-PPC/CCSPC/RCCN and Search-G components, governance, status,
coordination, and cited experiment/result artifacts.  Repository searches
covered decision fibers, future challenge, response signatures, source
removal, reread, recomputation, and staticization.  Primary literature was
checked before forming the theorem.

\section{DUPLICATION / CONTRADICTION AUDIT}

REREAD1 already proves that explicit static causal cascades can be streamed by
rereading.  REACQ1 proves deterministic selector and trace reacquisition is
polynomial.  STATICNP1 proves total NP-level continuation search is polynomial
under $P=NP$.  ONETRANS1 and FORKFACTOR1 close fixed replay and
wrapper-covariant internal forks.  TESTS-0019 closes bit-local pointer
persistence.  None of those reports states the exact allowed-response
intersection equivalence (2).  FIBERDICH1 does not re-claim their results; it
uses them to classify the two sides of the new exact boundary.  No verified
claim is contradicted.

\section{WHAT CHANGED RELATIVE TO ROOT MAIN}

Root MAIN correctly leaves the pointed source compiler and forbidden endpoint
\OPEN.  Its latest executive synthesis predates TESTS-0019 and does not state
the answer-fiber criterion.  The LSPV-4 region is sharpened to require every
online proposal to say whether its responses factor through SAT-answer
fibers.  Answer-factorable relations cannot force superconstant retained
source state; nonfactorable relations are source-specific and require a new
compiler.  This is a guard on the open target, not a proof that no such
compiler exists.

\section{ADOPTED / REJECTED / RETIRED PRIOR COMPONENTS}

Adopted: the HEC response-signature lower bound in its no-reread model,
REREAD1, REACQ1, STATICNP1, MODELBRIDGE-0020 pointer erasure, TESTS-0019
bit-local erasure, wrapper covariance, one-bit output, and full
SAME\_REALIZATION accounting.  Rejected: treating source deletion itself as
a static-SAT transfer; treating a response distinction inside one SAT-answer
fiber as information revealed by the SAT bit.  Retired:
\texttt{ANSWER-DERIVED-ONLINE-STATE1} as an unqualified superconstant
obligation.  Retained \OPEN: a source-fiber-separating full-language compiler
and one model-matched unconditional endpoint.

\section{STALE OR OVERSTATED CLAIMS FOUND}

Any wording that says a future challenge forces an arbitrary SAT decider to
retain the response signature is overstated unless the response relation,
source partition, and static-to-online compiler are explicit.  The HEC core
itself states the restricted-model scope and is not stale.  Root MAIN is
conservative.  The stale reading, not the restricted theorem, is retired.

\section{Status and next handoff}

\begin{center}\small
\begin{tabular}{p{0.47\linewidth}p{0.20\linewidth}p{0.25\linewidth}}
\toprule
Claim & Status & Scope\\ \midrule
\texttt{DECISION-FIBER-ONLINE-DICHOTOMY1} & \EXACT & Deterministic total response relation; review pending\\
Finite relation/hostile audit & \COMP & 6,561 relations; zero failures\\
NP-level response search under $P=NP$ & \COND & Balanced total polynomial verifier\\
\texttt{ANSWER-DERIVED-ONLINE-STATE1} & \RETIRED & Superconstant state claim\\
\texttt{SOURCE-FIBER-SEPARATING-COMPILER1} & \OPEN & Static-to-online transfer and all costs\\
Model-matched forbidden endpoint & \OPEN & Separation-strength\\
P versus NP & \OPEN & Arbitrary uniform computation\\
\bottomrule
\end{tabular}
\end{center}

The precise synthesis handoff is: before testing another pointer, fork, or
future challenge, compute the intersections in (2).  If they are nonempty,
publish the one-bit responder.  If any is empty, identify the same-answer
source pair and prove why every ordinary static SAT decider, rather than only
the added protocol, must preserve their distinction in one charged
realization.  Then audit that exact compiler against reread, rerun, semantic
and algebraic oracles, compact branching programs, alternate encodings, and
the easy families.

\begin{thebibliography}{9}
\bibitem{Yao79} A. C.-C. Yao, ``Some complexity questions related to
distributive computing,'' \emph{STOC 1979}, 209--213.
\url{https://doi.org/10.1145/800135.804414}
\bibitem{BKM18} L. Boczkowski, I. Kerenidis, and F. Magniez,
``Streaming communication protocols,'' \emph{ACM TOCT} 10(4), 2018.
\url{https://doi.org/10.1145/3276748}
\bibitem{BGS75} T. Baker, J. Gill, and R. Solovay,
``Relativizations of the $P=?NP$ question,'' \emph{SIAM J. Comput.} 4(4), 1975.
\url{https://doi.org/10.1137/0204037}
\bibitem{AW09} S. Aaronson and A. Wigderson,
``Algebrization: a new barrier in complexity theory,'' \emph{ACM TOCT} 1(1), 2009.
\url{https://doi.org/10.1145/1490270.1490272}
\bibitem{RR97} A. Razborov and S. Rudich, ``Natural proofs,''
\emph{J. Comput. Syst. Sci.} 55(1), 1997.
\url{https://doi.org/10.1006/jcss.1997.1494}
\bibitem{FLMV05} L. Fortnow, R. Lipton, D. van Melkebeek, and A. Viglas,
``Time-space lower bounds for satisfiability,'' \emph{J. ACM} 52(6), 2005.
\url{https://doi.org/10.1145/1101821.1101828}
\end{thebibliography}

% END_OF_VERIFIED_REPORT::LSPV-4::FIBERDICH1
\end{document}
